\documentclass[11pt,a4paper]{article}

% Pakete
\usepackage[utf8]{inputenc}
\usepackage[T1]{fontenc}
\usepackage[ngerman]{babel}
\usepackage[left=18mm,right=18mm,top=18mm,bottom=18mm]{geometry}
\usepackage{helvet}
\renewcommand{\familydefault}{\sfdefault}
\usepackage{amsmath,amssymb}
\usepackage{graphicx}
\usepackage{tcolorbox}
\tcbuselibrary{skins,breakable}
\usepackage{enumitem}
\usepackage{tabularx}
\usepackage{colortbl}
\usepackage{xcolor}
\usepackage[hidelinks]{hyperref}

% Fachfarbe Physik
\definecolor{physik}{HTML}{2c5aa0}
\definecolor{physiklight}{HTML}{e8f0fa}
\definecolor{loesung}{HTML}{c62828}

% Box-Stile
\tcbset{
    formelbox/.style={
        colback=white,
        colframe=physik,
        boxrule=2pt,
        arc=2pt,
        left=8pt,
        right=8pt,
        top=8pt,
        bottom=8pt
    },
    loesungbox/.style={
        colback=red!5,
        colframe=loesung,
        boxrule=1pt,
        arc=2pt,
        left=6pt,
        right=6pt,
        top=6pt,
        bottom=6pt
    }
}

% Kopfzeile
\usepackage{fancyhdr}
\pagestyle{fancy}
\fancyhf{}
\lhead{\small Physik 12 GK}
\chead{\small \textcolor{loesung}{\textbf{MUSTERLÖSUNG}}}
\rhead{\small Bohr: Rechnen}
\renewcommand{\headrulewidth}{0.4pt}

% Keine Einrückung
\setlength{\parindent}{0pt}
\setlength{\parskip}{6pt}

\begin{document}

% Titel
{\Large\textbf{Bohr-Atommodell: Rechnen}} \hfill \textcolor{loesung}{\textbf{MUSTERLÖSUNG}}\\[0.3em]
{\normalsize Fr 30.01.2026 | Physik 12 GK}

\vspace{0.3em}
\hrule height 1pt
\vspace{0.6em}

% ============================================================
% AUFGABE 1
% ============================================================

\textbf{Aufgabe 1: Energieniveaus berechnen} \hfill (3 BE)

\begin{center}
\begin{tabular}{|c|c|c|c|c|c|}
\hline
$n$ & 1 & 2 & 3 & 4 & $\infty$ \\
\hline
$E_n$ (eV) & \textcolor{loesung}{\textbf{-13,6}} & \textcolor{loesung}{\textbf{-3,4}} & \textcolor{loesung}{\textbf{-1,51}} & \textcolor{loesung}{\textbf{-0,85}} & 0 \\
\hline
\end{tabular}
\end{center}

\begin{tcolorbox}[loesungbox]
$E_1 = -13{,}6/1 = \mathbf{-13{,}6}$ eV \quad
$E_2 = -13{,}6/4 = \mathbf{-3{,}4}$ eV \quad
$E_3 = -13{,}6/9 = \mathbf{-1{,}51}$ eV \quad
$E_4 = -13{,}6/16 = \mathbf{-0{,}85}$ eV
\end{tcolorbox}

% ============================================================
% AUFGABE 2
% ============================================================

\textbf{Aufgabe 2: Emission -- Übergang $n = 3 \to n = 2$} \hfill (4 BE)

a) Berechne die Energie des emittierten Photons.

\begin{tcolorbox}[loesungbox]
$E_{\text{Photon}} = |E_3 - E_2| = |(-1{,}51) - (-3{,}4)| = |-1{,}51 + 3{,}4| = \mathbf{1{,}89 \text{ eV}}$
\end{tcolorbox}

b) Berechne die Wellenlänge und gib die Farbe an.

\begin{tcolorbox}[loesungbox]
$\lambda = \dfrac{h \cdot c}{E} = \dfrac{1240 \text{ eV} \cdot \text{nm}}{1{,}89 \text{ eV}} = \mathbf{656 \text{ nm}}$ \quad (H$_\alpha$-Linie, \textbf{rot})
\end{tcolorbox}

% ============================================================
% AUFGABE 3
% ============================================================

\textbf{Aufgabe 3: Rydberg-Formel anwenden} \hfill (3 BE)

Übergang $n = 4 \to n = 2$:

\begin{tcolorbox}[loesungbox]
\begin{align*}
\frac{1}{\lambda} &= R_H \cdot \left( \frac{1}{n_1^2} - \frac{1}{n_2^2} \right) = 1{,}097 \cdot 10^7 \text{ m}^{-1} \cdot \left( \frac{1}{4} - \frac{1}{16} \right) \\[0.5em]
&= 1{,}097 \cdot 10^7 \cdot \frac{3}{16} = 2{,}057 \cdot 10^6 \text{ m}^{-1} \\[0.5em]
\lambda &= \frac{1}{2{,}057 \cdot 10^6} \text{ m} = 4{,}86 \cdot 10^{-7} \text{ m} = \mathbf{486 \text{ nm}} \quad \text{(H$_\beta$-Linie, blau-grün)}
\end{align*}
\end{tcolorbox}

% ============================================================
% AUFGABE 4
% ============================================================

\textbf{Aufgabe 4: Absorption} \hfill (3 BE)

a) Energie für Übergang $n = 1 \to n = 3$:

\begin{tcolorbox}[loesungbox]
$E_{\text{Photon}} = |E_1 - E_3| = |(-13{,}6) - (-1{,}51)| = \mathbf{12{,}09 \text{ eV}}$
\end{tcolorbox}

b) Wellenlänge und Spektralbereich:

\begin{tcolorbox}[loesungbox]
$\lambda = \dfrac{1240}{12{,}09} = \mathbf{103 \text{ nm}}$ \quad (\textbf{UV-Strahlung}, Lyman-Serie)
\end{tcolorbox}

% ============================================================
% AUFGABE 5 (AFB III)
% ============================================================

\textbf{Aufgabe 5: Transfer -- Helium-Ion He$^+$} \hfill (3 BE)

\begin{tcolorbox}[loesungbox]
Für He$^+$ gilt: $E_n = \dfrac{-13{,}6 \cdot Z^2}{n^2} = \dfrac{-13{,}6 \cdot 4}{n^2}$

\begin{align*}
E_3(\text{He}^+) &= \frac{-13{,}6 \cdot 4}{9} = \mathbf{-6{,}04 \text{ eV}} \\[0.3em]
E_2(\text{He}^+) &= \frac{-13{,}6 \cdot 4}{4} = \mathbf{-13{,}6 \text{ eV}} \\[0.3em]
E_{\text{Photon}} &= |E_3 - E_2| = |-6{,}04 - (-13{,}6)| = \mathbf{7{,}56 \text{ eV}} \\[0.3em]
\lambda &= \frac{1240}{7{,}56} = \mathbf{164 \text{ nm}} \quad \text{(\textbf{UV-Strahlung})}
\end{align*}
\end{tcolorbox}

% ============================================================
% AUFGABE 6 (AFB III)
% ============================================================

\textbf{Aufgabe 6: Grenzen des Bohr-Modells} \hfill (2 BE)

\begin{tcolorbox}[loesungbox]
\textbf{Warum versagt das Modell bei Helium?}
\begin{itemize}[noitemsep,topsep=2pt]
\item Das Bohr-Modell berücksichtigt nur die \textbf{Coulomb-Wechselwirkung} zwischen Kern und \emph{einem} Elektron.
\item Bei zwei Elektronen gibt es zusätzlich die \textbf{Elektron-Elektron-Abstoßung}, die nicht erfasst wird.
\end{itemize}

\textbf{Trotzdem sinnvoll?}
\begin{itemize}[noitemsep,topsep=2pt]
\item \textbf{Ja:} Das Modell erklärt anschaulich die Quantisierung und diskrete Spektren.
\item Es liefert korrekte Ergebnisse für wasserstoffähnliche Systeme (H, He$^+$, Li$^{2+}$).
\item Es fördert \textbf{Modellkompetenz}: Schüler lernen, dass Modelle Grenzen haben.
\end{itemize}
\end{tcolorbox}

% ============================================================
% AUFGABE 7
% ============================================================

\textbf{Aufgabe 7: Spektralanalyse} \hfill (2 BE)

a) Elemente in der Sonnenatmosphäre:

\begin{tcolorbox}[loesungbox]
\textbf{Wasserstoff (H), Helium (He), Natrium (Na), Calcium (Ca)}

(Dies sind die in der Simulation nachweisbaren Elemente. In Wirklichkeit wurden über 60 Elemente im Sonnenspektrum gefunden.)
\end{tcolorbox}

b) Warum Absorptionslinien?

\begin{tcolorbox}[loesungbox]
Die \textbf{heiße Sonnenoberfläche} (ca. 5500°C) strahlt ein \textbf{kontinuierliches Spektrum} aus (Wärmestrahlung).

Dieses Licht durchquert die \textbf{kühlere Sonnenatmosphäre}. Dort werden Photonen mit genau passender Energie von Atomen \textbf{absorbiert} $\to$ Elektronen springen auf höhere Bahnen.

$\Rightarrow$ Im Spektrum fehlen diese Wellenlängen $\to$ \textbf{dunkle Linien} (Fraunhofer-Linien).
\end{tcolorbox}

\vfill

\begin{center}
\rule{10cm}{0.5pt}

\textbf{Gesamt:} \textcolor{loesung}{\textbf{20}} / 20 BE
\end{center}

\end{document}
